\documentclass[border=0pt]{standalone}
% The shared preamble of every docs diagram (docs/diagrams/<name>.tex).
% dvisvgm's driver first, so patterns and opacity survive into the SVG.
\def\pgfsysdriver{pgfsys-dvisvgm.def}
\usepackage[siunitx,RPvoltages]{circuitikz}
\usepackage{xcolor}
% Latin Modern: outline fonts for every glyph a label uses, µ included.
\usepackage[T1]{fontenc}
\usepackage{lmodern}
\usetikzlibrary{arrows.meta,backgrounds,calc,patterns,positioning}
\tikzset{
  % A white card, so diagrams read in mdBook's dark themes too.
  every picture/.append style={show background rectangle,
    background rectangle/.style={fill=white,inner sep=4pt}},
  % Layout geometry: layers told apart by fill and hatching, never colour alone.
  layer a/.style={draw=black,line width=0.6pt,fill=gray!20},
  layer b/.style={draw=black,line width=0.6pt,fill=gray!45},
  layer c/.style={draw=black,line width=0.6pt,fill=white,postaction={pattern=north east lines}},
  layer d/.style={draw=black,line width=0.6pt,fill=white,postaction={pattern=crosshatch dots}},
  marker/.style={draw=black,dashed,line width=0.6pt},
  % A measured distance and its label.
  dim/.style={{Latex[length=2mm]}-{Latex[length=2mm]},line width=0.5pt},
  dim label/.style={font=\small,fill=white,inner sep=1pt},
  % Where a rule fails: thick, dashed and labelled, not only red.
  violation/.style={draw=red!75!black,line width=1.4pt,dash pattern=on 3pt off 2pt},
  violation label/.style={font=\small\bfseries,text=red!75!black},
  caption/.style={font=\small},
}

\begin{document}
\begin{circuitikz}
  % Cross-section: a met1 wire over the substrate, plate and fringe fields.
  \fill[layer a] (0,0) rectangle (6,-0.25);
  \draw[line width=0.6pt] (-0.3,0) -- (6.3,0);
  \node[caption,anchor=west] at (6.4,-0.12) {substrate (ground)};
  \fill[layer b] (2.5,1.4) rectangle (3.5,1.75);
  \node[caption] at (3,1.575) {a};
  \node[caption,anchor=west] at (3.6,1.9) {met1, \SI{0.14}{\micro\metre} wide};
  \draw[dim] (5.6,0) -- (5.6,1.4) node[dim label,midway,right=2pt] {\SI{1.376}{\micro\metre}};
  % Plate field straight down, fringe field from the sides.
  \foreach \x in {2.7,3.0,3.3} \draw[-{Latex[length=1.5mm]}] (\x,1.35) -- (\x,0.05);
  \draw[-{Latex[length=1.5mm]}] (2.45,1.55) to[out=180,in=90] (1.5,0.05);
  \draw[-{Latex[length=1.5mm]}] (3.55,1.55) to[out=0,in=90] (4.5,0.05);
  \node[caption] at (1.2,1.0) {fringe};
  \node[caption] at (3.0,0.65) [fill=white,inner sep=1pt] {plate};
  % The network GPurify extracts from the 10 um wire.
  \draw (0,-1.4) node[circ,label=above:{a}]{}
    to[R, l=\SI{8.93}{\ohm}] (3,-1.4) node[circ,label=above:{a:1}]{};
  \draw (0,-1.4) to[C, l_=\SI{0.859}{\femto\farad}] (0,-3) node[ground]{};
  \node[caption,anchor=west] at (3.6,-2.2) {$\SI{25.78}{aF/\micro\metre^2}\times\SI{1.4}{\micro\metre^2}$};
  \node[caption,anchor=west] at (3.6,-2.7) {$+\ \SI{40.57}{aF/\micro\metre}\times\SI{20.28}{\micro\metre}$};
\end{circuitikz}
\end{document}
